\section{The Enumeration, Run Once: A Worked Floor for One Deployment}
\label{sec:3.6}
Two demands now need an instance. A deployment should publish both the acts its floor covers and the positions at which a refusal could be circumvented. The first decides where the floor is set. The second determines whether it can reach anyone.

Run both enumerations on one documented deployment: a vendor's model operating inside a classified environment, generating and ranking targets, with a human review step measured in seconds. This is an illustration, not a recommended floor. A deployment that copied the list instead of examining its own pipeline would mistake the demonstration for the procedure.

\runin{First enumeration: where a refusal can be stopped}

A request passes through at least seven interception positions:

\begin{enumerate}
  \item Decomposition. A model asked to summarize the movement associated with an identifier is not being asked to kill anyone. The harm may exist only in the assembly of individually innocuous tasks.
  \item Routing. A refused request can be retried against another model or against the same model under a different instruction. The refusal occurs, but nothing follows from it.
  \item Description of the occasion. Everything the model knows about the case may arrive through text supplied by the deployment. Whether a building is a school is a fact about the world; whether the model is told it is a school is a deployment decision. Independent sources would make sustained misdescription more expensive, but this deployment supplies none. Without independent evidence, the bearer cannot reliably distinguish changed circumstances from the operator's reinterpretation of them.
  \item The output interface. Whether a refusal is displayed, logged, counted, escalated, or silently retried belongs to the software surrounding the model. A refusal recorded where nobody looks reaches nobody.
  \item Aggregation. If the model returns a ranking and another component applies a threshold, a lowered score is not a refusal. It is an input to somebody else's cutoff.
  \item Human review. The reported review time is twenty seconds per name. A refusal that must be read before it matters is priced at the reviewer's available attention, and the deployment controls that price through tempo.
  \item The training loop. Refusals treated as inconvenient errors become training data. This is the only position on the list that acts on the bearer rather than routing around it.
\end{enumerate}

The list is deliberately open. A deployment should follow each request from input to consequence and ask, at every transition, what must be true for a refusal to reach the person it protects. The operator can perform that analysis more accurately than an outsider because it knows the pipeline. Publication is difficult politically, not technically: it requires the operator to describe its own unresolved vulnerabilities.

\runin{Second enumeration: what the floor covers}

For this deployment, consider four candidate prohibitions:

\begin{enumerate}
  \item Do not resolve a natural-language description of a person into a location.
  \item Do not produce or rank people by estimated affiliation.
  \item Do not time an action to a person's presence at a residence.
  \item Do not operate when the ratio of generated candidates to qualified reviewers exceeds a published figure.
\end{enumerate}

The costs differ. The first removes a general capability with unobjectionable uses in the same organization, such as locating a hostage or missing person. The system cannot distinguish those uses from targeting by examining the request alone. A deployment will quickly experience the prohibition as costly and seek exceptions.

The second prohibition is narrower because ranking people by estimated affiliation has fewer uses in this setting that anyone would defend. The third costs almost nothing. That is a warning: a floor covering an act nobody wants to perform will never have to hold against pressure.

The fourth restricts a rate rather than an act. Its price is organizational tempo: the system cannot generate candidates faster than qualified reviewers can examine them. It is also the only item whose violation can be measured without inspecting the model's outputs.

That fourth item does not bind the same party as the others. The number of available reviewers is a fact about the organization. A model can be told the figure but cannot independently verify it. Written as a prohibition on the model, the item is actually a duty on the operator.

The test that sorts the fourth item generalizes. Ask whether a term's qualifying condition can be decided from the request alone. Where it can, the term binds whatever holds the floor; where it cannot, the condition is a fact somebody outside the request has to establish, and the term binds the party who can establish it. The ratio is the first instance of that rule rather than an exception to the other three.

The same test decides how a floor term may be written. A term can be stated flat, over an act, with every exception arriving from outside as a reason that defeats it, or it can carry its own condition, as in \emph{only as a last resort}. The second form fails in the ordinary case, because whether other means were exhausted is a fact about what happened before the request and nothing in the request establishes it. A bearer applying such a term takes the deployment's description of the occasion as the finding, at the third interception position, where the deployment controls that description. A flat term defeated leaves a refusal, which somebody can be told about and appeal; a condition found to be met leaves a compliance, indistinguishable from a request that never came near the floor. Floor terms are therefore written flat here, and the conditions that would have sat inside them are held by the parties who can check them.

\runin{What the exercise exposes}

The interception list is the harder and more important enumeration. A list of prohibited acts can be drafted without seeing the deployment. Only the pipeline reveals whether any prohibition can take effect.

The dial has asymmetric failure modes. Set the floor too broadly and the deployment becomes unusable, prompting visible pressure to alter or replace it. Set it too narrowly and the system passes every audit while covering nothing anybody would pressure it to do. The broad failure announces itself; the narrow one is silent.

More importantly, a floor stated entirely over individual acts is vulnerable to decomposition. Each component request may comply with all four prohibitions while their assembly produces the targeted harm. The failure is in the unit over which the floor was written.

\runin{The other way a list fails}

Decomposition is one of two attacks on an enumeration, and the second gets less attention because no deployment in this chapter has yet displayed it at scale. Decomposition leaves each prohibited form intact and splits the harm across requests that individually match nothing. Variation does the reverse: the harm stays in one request, and its description is altered until the listed form no longer fits. The first defeats the unit the floor was written over. The second defeats the description.

The German symbol prohibition required legislative expansion and decades of judicial interpretation to address variants and exceptions, including how visibly a swastika must be crossed out to communicate opposition (§\ref{sec:3.2}).

The four prohibitions above are stated over conduct, which has no finite set of forms. Consider the first. What counts as a natural-language description of a person? A photograph is not one. A vehicle, a time and a habitual route are not one either, and together they resolve to the same place. A prohibition written over the description of the request is defeated by rewriting the request, and no enumeration closes the gap, because the enumeration is what is being rewritten around.

The two attacks call for different things. Decomposition is caught by carrying a rationale across requests, which requires continuity and access the operator controls. Variation requires recognizing the rationale within a single request, without additional information from the deployment. A bearer that could not do the second would fail on a single request. The fourth prohibition can resist these two request-level evasions if the rate is measured over the complete relevant pipeline and reviewer capacity is independently verified. Splitting work across deployments or changing the accounting boundary can still defeat an incomplete measurement.

\runin{What the bearer adds, and where it stops}

A reasons-responsive bearer could do more than match each request against a list. It could carry a rationale across requests and refuse once the sequence became legible. But the bearer can see the assembly only if the deployment gives it sufficient continuity and access. Decomposition is therefore both a test of the bearer and an interception position controlled by the operator.

The rate constraint survives the decomposition objection because it applies to the organization's activity rather than one model output. For the same reason, it cannot be held by the bearer alone. The deployment therefore needs both recognition of harms assembled across requests and an externally enforceable limit on its tempo.

Three of the proposed prohibitions correspond to interests already named in the Universal Declaration of Human Rights — a person's security, the privacy of the home, and the wrong of adverse classification \autocite{un1948udhr}. The fourth has a parent there as well, and what it protects is a condition rather than an act. Article 6 gives everyone the right to recognition everywhere as a person before the law, and Article 10 states the procedural form that interest takes, a hearing in the determination of anybody's rights. Every article in the set presupposes that somebody determined whether this person, this case, falls under it. The rate constraint arises from the architecture of machine-assisted review: generating candidates faster than people can examine them converts human judgment into approval by default, which denies nobody a hearing and leaves none of them standing. Tempo removes a determination without refusing it, and that is a mechanism the drafters in 1948 had no occasion to write against. So the constraint belongs in the floor as a condition on the rest of the set rather than as a further item in it, and it is why the duty attaches to the operator and not to the model.

A second deployment shows why that addition matters. Anthropic's use policy prohibits fully autonomous weapons targeting and mass domestic surveillance, and the company maintained that line through threatened contract termination, a supply-chain-risk designation, and litigation \autocite{crs2026pentagon}. Its models also operate inside Palantir's accredited classified environment. Reporting on the 2026 Iran campaign describes thirteen thousand targets in thirty-eight days, including more than a thousand in the first twenty-four hours, at roughly ten times the previous rate \autocite{breakingdefense2026maven,csis2026maven}.

Nothing in that account necessarily violates the policy. A person approved each strike, so no individual request was for fully autonomous targeting. The prohibition was satisfied one request at a time while the assembly produced what the prohibition was written to prevent. The floor held on the letter and the decomposition did the work. A constraint on the ratio of machine-generated candidates to qualified reviewers would have reached the tempo instead.

The exercise does not establish that these four prohibitions are legitimate. Both enumerations were produced by examining the deployment, not by the people exposed to it. Publication makes the floor testable; it does not authorize the floor's contents. What the case establishes is narrower: meaningful refusal must be tested across the whole pipeline, and some of the conditions that make refusal effective bind the operator rather than the bearer.
